\section{Introduction}
\label{sec:intro}
Real-time conversation rewards two behaviors that pull in opposite directions.
For a greeting or an easy factual question, a speech agent should answer as
soon as it understands the turn. For a difficult knowledge or reasoning
question, speaking immediately can lock the interaction into a fluent but
incorrect answer; taking time to reason or ask a stronger model is preferable.
The choice is therefore whether to answer now or wait for stronger reasoning.

Native duplex models already learn a fine-grained listen/speak policy
\citep{defossez2024moshi,openbmb2026minicpmo}. Recent systems extend this idea
with tokenized control. BayLing-Duplex uses four dialogue-state tokens for
turn-taking and interruption \citep{bayling2026duplex}; DuplexOmni uses a
\texttt{[THINK]} token to request an asynchronous reasoning layer
\citep{huang2026duplexomni}; and Listen-Write-Speak encodes concurrent
listening, writing, and speaking with a token schema
\citep{zhang2026liberating}. These systems show that special-token control is
expressive. They also make the routing rule part of autoregressive behavior:
the model must learn rare control events together with speech content,
timing, and interruption. BayLing-Duplex, for example, reweights rare state
tokens, adds timing-focused preference optimization, and retains an auxiliary
language-modeling loss. DuplexOmni and Listen-Write-Speak construct large
temporally annotated training corpora. Updating the rule consequently means
updating or retraining the interaction model.

Turn-taking supervision also does not directly answer the competence question.
A model may be ready to take the floor while lacking the knowledge needed for
a correct response. Conversely, routinely invoking a strong expert wastes cost
and makes easy turns wait. What is needed is a failure-risk estimate that is
available at the native speak boundary, does not disturb the learned duplex
policy, and can be recalibrated as the desired accuracy--latency operating point
changes.

We investigate whether this decision can be recovered from a frozen duplex
model's own representations. Hidden-state probes can predict correctness in
text models \citep{azaria2023internal,kossen2024semantic,lugoloobi2026failures},
but a real-time speech setting introduces three shifts. The input is audio
rather than text; duplex fine-tuning reshapes representations for interaction;
and the read must be available at the listen-to-speak transition, before a
complete answer can be inspected. A successful gate must survive all three and
produce a useful serving policy, not merely an offline AUC measured after the
answer has already been generated.

\begin{figure}[t]
  \centering
  \includegraphics[width=\linewidth]{teaser_v2}
  \caption{\textbf{Post-training routing at the native speak boundary.}
  \textbf{(1)} Speech streams in 1\,s chunks to a frozen full-duplex local
  model. \textbf{(2)} When the model commits to speak, an L22 linear readout
  scores local failure risk in tens of milliseconds; the triggering call may
  already contain the first text unit. \textbf{(3)} Low-risk turns continue
  immediately, while high-risk turns are sent to a stronger expert and the
  talker holds the floor. \textbf{(4)} The expert answer is voiced by the local
  talker. On all 240 internal queries, the aggressive policy raises accuracy
  from 40.8\% to 62.9\% at 51.7\% escalation.}
  \label{fig:hero}
\end{figure}

Our system freezes MiniCPM-o~4.5 \citep{openbmb2026minicpmo} and trains only a
regularized logistic readout on cached answer-onset states. This separates
\emph{interaction control}---when the model takes the floor---from
\emph{competence control}---whether its answer should be trusted. The readout
adds no language-model forward pass, leaves local responses unchanged, and
turns the expert-call budget into a threshold that can be adjusted after
training. On the expert path, the question is sent to a stronger model and the
returned answer is spoken by the local talker.

Our approach treats competence as a post-training serving decision. We
implement think-or-talk routing at the native speak boundary with a linear gate,
without adding control tokens, updating the backbone, or repeating duplex
training. We trace the readable signal through duplex and modality shifts:
text-trained probes transfer to unseen human speech, while the matched native
ablation identifies answer-onset aggregation as important for robust readout.
We evaluate the gate as a complete serving policy rather than only an offline
classifier. On the 240-query internal set, selective handoff improves accuracy
by 22.1 points while reducing mean server waiting time by 35.3\% relative to
always escalating. Across four external speech-QA pools, the aggressive policy
improves macro accuracy by 21.4 points over local answering and by 5.9 points
over a matched random router at the same expert-call rate in every external
evaluation pool.

The result is a practical division of labor: the duplex model preserves the
fast interaction for which it was trained, while a separately calibrated gate
decides when additional reasoning is worth the delay.
